% Part: sets-functions-relations
% Chapter: functions
% Section: functions-relations

\documentclass[../../../include/open-logic-section]{subfiles}


\begin{document}

\olfileid[ar]{sfr}{fun}{rel}

\olsection{الدوال بوصفها علاقات}

\begin{explain}
الدالة التي تقابل !!{element}s~$A$ بـ!!{element}s~$B$ تعيّن علاقة
بين $A$ و~$B$: وهي التي تقوم بين $x$ و$y$ إذا وفقط إذا كان
$f(x) = y$. وقد نريد، اختزالًا لأصول البناء الرياضي مثلًا،
\emph{مماهاة}~$f$ بهذه العلاقة، أي بمجموعة الأزواج.
فأي العلاقات يجوز أن تؤخذ دوال على هذا الوجه؟
\end{explain}

\begin{defn}[الرسم البياني لدالة]
للدالة $f\colon A \to B$ \emph{رسم بياني} للدالة~$f$، وهو العلاقة
$R_f \subseteq A \times B$ ذات التعريف
\[
R_f = \Setabs{\tuple{x,y}}{f(x) = y}.
\]
\end{defn}

\begin{explain}
تقضي الامتدادية بوحدة رسم الدالة البياني. بل إن امتدادية المجموعات
تسوّغ مباشرة امتدادية الدوال التي استعملناها ضمنًا: متى اتحد
مجال $f$ و~$g$ ومجالهما المقابل، واتفقت قيمهما كلها، كانتا واحدة.

وبالعكس، كل علاقة تستوفي شرط «الدالية» هي رسم بياني لدالة.
\end{explain}

\begin{prop}\ollabel{prop:graph-function}
افرض أن $R \subseteq A \times B$ تحقق الشرطين:
\begin{enumerate}
\item اجتماع $Rxy$ و$Rxz$ يقتضي $y = z$؛ و
\item يوجد لكل $x \in A$ عنصر $y \in B$ يحقق $\tuple{x,
y} \in R$.
\end{enumerate}
فتكون $R$ رسم الدالة $f\colon A \to B$ التي نضعها بالشرط:
$f(x) = y$ إذا وفقط إذا كان $Rxy$.
\end{prop}

\begin{proof}
إذا وجد $y$ يحقق $Rxy$ امتنع أن يوجد $z \neq
y$ يحقق $Rxz$، وإلا خولف شرط~$R$. فمتى وجد $y$ مع $Rxy$ كان
هذا $y$ واحدًا؛ ومع ضمان الوجود تكون $f$ سليمة التعريف.
ومن تعريفها يظهر $R_f = R$.
\end{proof}

\begin{explain}
فكل دالة $f\colon A \to B$ لها رسم، أي علاقة بين $A$ و~$B$
محتواة في~$A \times B$ يضبطها $f(x) = y$. وكل علاقة~$R
\subseteq A \times B$ تحقق شرطي \olref{prop:graph-function} رسم
لدالة~$f \colon A \to B$. ولشدة هذه الصلة يجوز أن نأخذ الدالة
رسمها نفسه، أي علاقة مخصوصة أو مجموعة مخصوصة من الصفوف المرتبة.
\oliflabeldef{sfr:rel:ref:sec}{والمقصود بالمماهاة هنا ما بُيّن في
\olref[sfr][rel][ref]{sec}: ليست دعوى ميتافيزيقية في ماهية الدوال،
بل إجازة نافعة لـ\emph{معاملتها} مجموعات معينة. ومن نفع ذلك أن }{}
العمليات التي أجريناها على العلاقات يمكن إجراء نظائرها على الدوال
(انظر \olref[sfr][rel][ops]{sec})؛ ومنها ما يأتي.
\end{explain}

\begin{defn}\ollabel{defn:funimage}
خذ دالة $f \colon A \to B$ وجزئية $C\subseteq A$.

\emph{تقييد}~$f$ على~$C$ هو الدالة~$\funrestrictionto{f}{C}\colon C \to B$
التي تحقق $(\funrestrictionto{f}{C})(x) = f(x)$ لكل $x \in C$؛ أي
$\funrestrictionto{f}{C} = \Setabs{\tuple{x, y} \in R_f}{x \in
C}$.

و\emph{تطبيق}~$f$ على~$C$ تعريفه $\funimage{f}{C} =
\Setabs{f(x)}{x \in C}$؛ ويسمى أيضًا \emph{صورة}~$C$ تحت~$f$.
\end{defn}

\begin{explain}
ومن التعريفين ينتج $\ran{f} =
\funimage{f}{\dom{f}}$ لكل دالة~$f$.
\oliflabeldef{sfr:rel:ops:sec}{وهذا نظير لما تقتضيه تعريفات
\olref[sfr][rel][ops]{sec} ومماهاة الدوال بالعلاقات، إلا أن تقييد
الدالة هنا إنما يقصر المدخل على~$C$ ولا يقصر المخرج على~$C$.
وأما المعكوسات والضروب النسبية فيحتاجان إلى زيادة تفصيل،
وسيأتي في \olref[inv]{sec} و\olref[cmp]{sec}.}{}
\end{explain}

\begin{explain}
\textbf{تنبيهان على لفظ الأصل.} سمى الأصل رسم~$f$ «علاقة على
$A\times B$»، مع أن اصطلاح الكتاب يجعل العلاقة على~$U$ جزءًا من
$U^2$؛ فقلنا علاقة بين $A$ و~$B$ محتواة في~$A\times B$. ثم إن
تقييد الدالة على~$C$ يخص الإحداثي الأول، فلا يطابق تقييد العلاقة
المتجانسة بالتقاطع مع~$C^2$ إلا على جهة النظير.
\end{explain}

\end{document}
